\documentclass[12pt]{article}
\usepackage[utf8]{inputenc}
\usepackage[T1]{fontenc}
\usepackage{hyperref}
\title{Bakung Kenyah Language Resources for K-12 Curriculum Development}
\author{Compiled by Claude Sonnet 5.5 (claude-sonnet-5-5) for Andr\'as Kornai}
\date{}
\begin{document}
\maketitle

\section{Introduction}
Bakung Kenyah (ISO 639-3 \texttt{boc}, Glottocode baku1265) belongs to the Kenyah group of Austronesian languages in Borneo. The family field in the prompt reads ``Bookkeeping'', which is evidently a data error, so no family claim is made from it. Search results indicate that Bakung is a name for a Kenyah dialect grouped with Lepo' Tepu', within the Mainstream Kenyah cluster (Usun Apau; ISO code xkl; roughly 50,000 speakers for the cluster as a whole) in North Kalimantan and Sarawak. I found no material specific to Bakung.

\section{Primary Educational Resources}
None found in the searches.

\section{Linguistic and Reference Materials}
\begin{itemize}
\item Mainstream Kenyah language: \url{https://en.wikipedia.org/wiki/Mainstream_Kenyah_language}
\item Kenyah languages: \url{https://en.wikipedia.org/wiki/Kenyah_languages}
\item Sellato (2009), classification of Bakung within Kenyah, Glottolog record: \url{https://glottolog.org/resource/reference/id/472374}
\end{itemize}

\section{Religious and Community Resources}
None found in the searches.

\section{Digital and Audio Resources}
\begin{itemize}
\item Omniglot page on Kenyah, which lists Bakung among its dialects: \url{https://www.omniglot.com/writing/kenyah.htm}
\end{itemize}

\section{Curriculum Development Context}
Whatever documentation exists is probably filed under the wider Kenyah label. A curriculum would need to establish what is specific to Bakung. Nothing found indicates school use.
\end{document}